\chapter{Normalization and the theory package}
\label{ch:ArchonFull600Theory_Package}

This chapter is Milestone 6 of the handoff: Lemmas 12, 13 and 15, and the
Theorem as the conjunction of its seven parts, each stated under its own
hypotheses. Lemma 14 is Lemma~\ref{lem:endpoint-transport}. Source line
numbers refer to \texttt{research/Full\_600cell\_Puzzle\_Theory.tex}.

\section{Lemma 12: geometric normalization}

\begin{lemma}[A symmetry fixing every facet]
  \label{lem:fix-every-cap}
  \lean{ArchonFull600Theory.eq_one_of_forall_smul_cap}
  \leanok
  \uses{def:symmetry-group, thm:face-counts}
  \textit{Source: Lemma 12, the trivial intersection.}
  A symmetry of \(P\) that maps every facet to itself is the identity.
\end{lemma}
\begin{proof}
  \leanok
  Four facets meet the facet \(\{0,30,31,32\}\) in one vertex each, so the
  symmetry fixes those four vertices, which form a basis of \(\mathbb H\).
  The source argues through the poles spanning \(\mathbb R^4\).
\end{proof}

\begin{theorem}[Lemma 12]
  \label{thm:lemma-twelve}
  \lean{ArchonFull600Theory.HypothesisE.conj_legalGroup,
    ArchonFull600Theory.HypothesisE.range_le_normalizer,
    ArchonFull600Theory.HypothesisE.legalGroup_inf_range,
    ArchonFull600Theory.HypothesisE.semidirectEquiv,
    ArchonFull600Theory.HypothesisE.existsUnique_mul,
    ArchonFull600Theory.HypothesisE.normalization}
  \leanok
  \uses{def:hypotheses-fce, thm:lemma-six, lem:fix-every-cap}
  % SOURCE: Lemma 12, lines 1267--1316
  % SOURCE QUOTE: Under R, F, and E, the geometric action of \(K\) normalizes \(G\), its
  % intersection with \(G\) is trivial, and \begin{equation} \langle G,K\rangle\cong G\rtimes K.
  % \end{equation}
  \textit{Source: Lemma 12.}
  Under R, F and E, with \(\kappa\colon K\to S_{\mathcal S}\) the geometric
  action: \(\kappa(k)G\kappa(k)^{-1}=G\) for every \(k\in K\),
  \(G\cap\kappa(K)=1\), and \(\langle G,\kappa(K)\rangle\cong G\rtimes K\)
  through \((g,k)\mapsto g\kappa(k)\), each element having a unique form
  \(g\kappa(k)\).
\end{theorem}
\begin{proof}
  \leanok
  E's field \texttt{geom\_capAction} conjugates each rotation of cap \(c\)
  into a rotation of cap \(k\cdot c\), and by F and Lemma 6 every cap
  rotation is legal. A legal element fixes every centre, and a
  \(\kappa(k)\) fixing every centre fixes every cap, so \(k=1\) by
  Lemma~\ref{lem:fix-every-cap}. The semidirect product maps injectively
  onto the join.
\end{proof}

\begin{lemma}[The legal group is larger than \(K\)]
  \label{lem:distinct-legal}
  \lean{ArchonFull600Theory.HypothesisC.card_symmetryGroup_lt}
  \leanok
  \uses{thm:selected-lemmas, thm:finite-symmetry-group}
  % SOURCE: Distinct symmetry and motion groups, lines 1357--1360
  % SOURCE QUOTE: Under C, the lower bound \(B>10^{151851}\) alone already separates the scale of
  % the legal action from the finite geometric group.
  \textit{Source: the remark on distinct groups.}
  Under R and C, \(|K|<|G|\).
\end{lemma}
\begin{proof}
  \leanok
  \(|K|=14400<10^{151851}<B\le|\Omega|\le|G|\).
\end{proof}

\section{Lemma 13: sound symmetry reduction}

\begin{lemma}[Lemma 13, the tetrahedral example]
  \label{lem:tetrahedral-conjugation}
  \lean{ArchonFull600Theory.tetraHalfTurn, ArchonFull600Theory.tetraThirdTurn,
    ArchonFull600Theory.tetraSymbols, ArchonFull600Theory.tetrahedralConjugation}
  \leanok
  \uses{def:word-distance}
  % SOURCE: Lemma 13, lines 1362--1427
  % SOURCE QUOTE: Normalization of \(G\) by \(K\) does not imply preservation of the chosen
  % \(1800\)-symbol word metric.
  % SOURCE QUOTE: Conjugating \(H\) by \(T\) gives a different half-turn.
  \textit{Source: Lemma 13, first statement.}
  In \(A_4\) with \(H=(1\,2)(3\,4)\) and \(T=(1\,2\,3)\), the conjugate
  \(THT^{-1}\) is a half-turn other than \(H\) and none of \(H\), \(T\),
  \(T^{-1}\). It is the word \(T^{-1},H,T\), and no shorter word gives it,
  while \(H\) is a word of one letter.
\end{lemma}
\begin{proof}
  \leanok
  Decision over \(A_4\), including every word of at most two letters.
\end{proof}

\begin{theorem}[Lemma 13, the quotient]
  \label{thm:symmetry-quotient}
  \lean{ArchonFull600Theory.PathCost, ArchonFull600Theory.quotientEdge,
    ArchonFull600Theory.quotientGoal, ArchonFull600Theory.PathCost.project,
    ArchonFull600Theory.PathCost.lift, ArchonFull600Theory.exists_goal_pathCost_iff,
    ArchonFull600Theory.goalCosts_eq, ArchonFull600Theory.isLeast_goalCosts_iff}
  \leanok
  % SOURCE: Lemma 13, lines 1371--1375
  % SOURCE QUOTE: A symmetry quotient preserves shortest-path distance to a goal set if the
  % symmetries act by cost-preserving automorphisms of the actual allowed state graph, preserve
  % the goal set
  \textit{Source: Lemma 13, second statement.}
  If a group \(L\) acts by cost-preserving automorphisms of a state graph
  and preserves the goal set, a state has a path of cost \(c\) to a goal
  exactly when its orbit has a path of cost \(c\) to a goal orbit, so the
  minimum costs agree.
\end{theorem}
\begin{proof}
  \leanok
  Paths project with their cost. An orbit edge represented by
  \(u\to v\) lifts at the state \(\ell u\) to the edge \(\ell u\to\ell v\),
  and the goal set is invariant.
\end{proof}

\section{Lemma 15: guarded setups}

\begin{theorem}[Lemma 15]
  \label{thm:guarded-setup}
  \lean{ArchonFull600Theory.frameEdge, ArchonFull600Theory.pathDistance,
    ArchonFull600Theory.setupLength, ArchonFull600Theory.withInverses,
    ArchonFull600Theory.pathDistance_le_of_embedding,
    ArchonFull600Theory.setupLength_withInverses_le,
    ArchonFull600Theory.pathDistance_withInverses_le, ArchonFull600Theory.withInverses_fixes,
    ArchonFull600Theory.reachable_withInverses_iff}
  \leanok
  \uses{thm:symmetry-quotient}
  % SOURCE: Lemma 15, lines 1470--1521
  % SOURCE QUOTE: enlarging the positive guarded alphabet by legal inverses cannot increase any
  % minimum unit-cost setup length.
  \textit{Source: Lemma 15.}
  On ordered target frames, adding the inverse letters cannot increase a
  minimum unit-cost setup length, nor a minimum weighted cost when the old
  letters keep their costs. Inverse letters fix the buffers when the
  letters do, and when every inverse transition has a path of positive
  letters, they do not enlarge the reachable set.
\end{theorem}
\begin{proof}
  \leanok
  Every old path is a new path of the same cost, and the least cost is an
  infimum in a complete lattice, with \(\top\) for unreachable frames.
\end{proof}

\begin{lemma}[Lemma 15 on the retained model]
  \label{lem:cap-guard}
  \lean{ArchonFull600Theory.HypothesisF.H_inv, ArchonFull600Theory.HypothesisF.T_inv,
    ArchonFull600Theory.capLetters, ArchonFull600Theory.capSetup_reachable_iff,
    ArchonFull600Theory.HypothesisR.CapSupport, ArchonFull600Theory.capGuard_fixes,
    ArchonFull600Theory.capGuardedSetup}
  \leanok
  \uses{thm:guarded-setup, def:hypotheses-fce, def:retained-model}
  % SOURCE: Lemma 15, lines 1494--1505
  % SOURCE QUOTE: In the retained cap guard, forbidding every cap in \begin{equation}
  % M(A)\cup M(B) \end{equation} is a sufficient condition for this fixed-buffer property.
  % SOURCE QUOTE: Moreover, \(H_c^{-1}=H_c\) and \(T_c^{-1}=T_c^2\), so inverse transitions
  % already have positive legal witnesses.
  \textit{Source: Lemma 15, the cap guard.}
  The inverse letters of a cap alphabet are legal and do not increase a
  setup length. Under F, \(H_c^{-1}=H_c\) and \(T_c^{-1}=T_c^2\), so they do
  not enlarge the reachable frames. Under the support property that a cap
  outside a piece's complete signature fixes its stickers, the source's
  ``rotates the cap and fixes its complement'', which R and F do not
  state, forbidding every cap of \(M(A)\cup M(B)\) makes every letter fix
  the buffer pieces \(A\) and \(B\).
\end{lemma}
\begin{proof}
  \leanok
  \(H_c\) and \(T_c\) are images of elements of order \(2\) and \(3\) of
  \(A_4\).
\end{proof}

\section{The theory package}

\begin{lemma}[Antipodal axes]
  \label{lem:antipodal-axes}
  \lean{ArchonFull600Theory.centralSymmetry, ArchonFull600Theory.antipode,
    ArchonFull600Theory.antipode_ne, ArchonFull600Theory.antipode_disjoint,
    ArchonFull600Theory.card_antipodalAxes, ArchonFull600Theory.antipodalAxes}
  \leanok
  \uses{def:cut-chambers, lem:metric-data, thm:binary-icosahedral}
  % SOURCE: Theorem, part 3, lines 1587--1590
  % SOURCE QUOTE: The 600 oriented poles describe 300 antipodal axes with two outer layers at
  % normalized distance \(121/125\)
  \textit{Source: the Theorem, part 3.}
  The central symmetry \(x\mapsto-x\) lies in \(K\). It maps each cap to a
  different cap with \(n_{-c}=-n_c\), so the caps form \(300\) antipodal
  pairs, and the two cap regions of a pair are disjoint.
\end{lemma}
\begin{proof}
  \leanok
  \(x\mapsto-x\) is \(L_{-1,1}\) with \(-1\in I^*\). A cap equal to its
  antipode would contain adjacent vertices \(v\) and \(-v\), at distance
  \(4\ne2/\phi\). The two cut functions of a pair cannot both be positive.
\end{proof}

\begin{lemma}[The groups are distinct]
  \label{lem:distinct-groups}
  \lean{ArchonFull600Theory.infinite_unitary, ArchonFull600Theory.infinite_specialOrthogonal,
    ArchonFull600Theory.card_capRotations}
  \leanok
  \uses{thm:double-cover, def:cut-chambers}
  % SOURCE: Theorem, part 2, lines 1580--1586
  % SOURCE QUOTE: These finite groups are distinct from \(SO(4)\), the rotation group of a cap and
  % the legal twist group.
  \textit{Source: the Theorem, part 2.}
  \(SO(4)\) is infinite, and the rotation group of a cap, the even
  permutations of its four vertices, has order \(12\).
\end{lemma}
\begin{proof}
  \leanok
  The unit quaternions \(t+\sqrt{1-t^2}\,i\) are infinitely many, and
  \(p\mapsto L_{p,1}\) is injective. \(4!/2=12\).
\end{proof}

\begin{lemma}[Inverse words]
  \label{lem:inverse-words}
  \lean{ArchonFull600Theory.exists_inverse_word, ArchonFull600Theory.inverseCapLetter,
    ArchonFull600Theory.HypothesisF.exists_inverse_word}
  \leanok
  \uses{lem:cap-words, lem:cap-guard}
  % SOURCE: Theorem, part 5, lines 1612--1622
  % SOURCE QUOTE: This is not a lower bound for each scramble; a recorded 1,000-letter word has
  % an inverse of at most 1,000 letters.
  \textit{Source: the Theorem, part 5.}
  Under F, every word in the \(1800\) letters \(H_c,T_c,T_c^{-1}\) has an
  inverse word of the same length.
\end{lemma}
\begin{proof}
  \leanok
  The reversed word of inverse letters, with \(H_c^{-1}=H_c\).
\end{proof}

\begin{theorem}[The Full 600-cell theory package]
  \label{thm:package}
  \lean{ArchonFull600Theory.TheoryPart1, ArchonFull600Theory.TheoryPart2,
    ArchonFull600Theory.TheoryPart3, ArchonFull600Theory.TheoryPart4,
    ArchonFull600Theory.TheoryPart5, ArchonFull600Theory.TheoryPart6,
    ArchonFull600Theory.TheoryPart7, ArchonFull600Theory.theoryPart1,
    ArchonFull600Theory.theoryPart2, ArchonFull600Theory.theoryPart3,
    ArchonFull600Theory.theoryPart4, ArchonFull600Theory.theoryPart5,
    ArchonFull600Theory.theoryPart6, ArchonFull600Theory.theoryPart7,
    ArchonFull600Theory.fullTheoryPackage}
  \leanok
  \uses{lem:coordinate-identification, lem:metric-data, thm:face-counts, lem:clique-face,
    thm:facet-completeness, thm:radial-homeomorph, thm:radial-topology,
    thm:binary-icosahedral, thm:finite-symmetry-group, thm:coxeter-presentation,
    thm:double-cover, lem:distinct-groups, lem:distinct-legal, lem:antipodal-axes,
    thm:signatures, lem:census-sums, lem:model-totals, thm:lemma-six, def:frame-wreath,
    thm:lemma-seven, lem:invariants, thm:invariant-bound, thm:selected-lemmas,
    lem:inverse-words, thm:lemma-twelve, lem:tetrahedral-conjugation,
    thm:symmetry-quotient, lem:endpoint-transport, thm:guarded-setup, lem:cap-guard,
    lem:scope-same-restriction, lem:scope-visible-disjoint, lem:disjoint-full-supports}
  % SOURCE: Theorem, lines 1557--1692
  % SOURCE QUOTE: The radius-two polytope \(P\) is the regular \(\{3,3,5\}\) 600-cell with edge
  % length \(2/\phi\), face vector \((120,720,1200,600)\), tetrahedral facets and icosahedral
  % vertex links.
  % SOURCE QUOTE: Under R and F, the legal group embeds in the stated product of
  % transported-frame wreath products
  % SOURCE QUOTE: The faithful geometric action on the complete retained subdivision normalizes
  % \(G\).
  % SOURCE QUOTE: Equality on a target orbit does not establish equality of full-state macros;
  % disjoint full supports suffice for commutation, while disjoint visible supports do not.
  \textit{Source: the Theorem.}
  The conjunction of seven parts, each under its own hypotheses. (1)
  Geometry and topology: unconditional for the coordinates, edge length,
  clique counts and faces, vertex link counts, five facets per edge and
  the radial homeomorphism; under E4 the facets are the \(600\)
  tetrahedra, under E3 the homology and simple connectivity, under E1 and
  E2 the identifications. (2) Symmetry: \(I^*\), \(K^+\) of order
  \(7200\), \(K\) of order \(14400\), the Coxeter relations and the
  presentation under E5, the double cover, and the distinct groups, the
  legal group under R and C. (3) The axes and cuts, the census, and under
  R the model counts, equivariance, signatures, addresses and ranks; under
  F local \(A_4\) generation. (4) The frame law; under R the embedding,
  the index \(2^{43}5^2\) and the sizes \(3600\), \(2400\), \(7200\) of the
  selected orbits; under I the upper bound; under C the alternating
  subgroup and the colour bound. (5) Under R and C the word metric, and
  under F inverse words. (6) Under R, F and E Lemma 12; the tetrahedral
  example and the quotient. (7) Lemma 14, Lemma 15 and the scope tests.
  Not formalized: the polar \(120\)-cell, flag regularity, the Schl\"afli
  symbol, and that an outer-layer turn fixes the middle layer.
\end{theorem}
\begin{proof}
  \leanok
  Each part assembles the lemmas of the earlier chapters and of this one.
\end{proof}
